\documentclass[11pt]{article}
\usepackage[utf8]{inputenc}
\usepackage[T1]{fontenc}
\usepackage{amsmath,amssymb}
\usepackage{graphicx}
\usepackage{booktabs}
\usepackage{array}
\usepackage[margin=1in]{geometry}
\usepackage{hyperref}
\hypersetup{colorlinks=true, linkcolor=blue, citecolor=blue}
\usepackage{natbib}
\bibliographystyle{unsrtnat}

\title{Distinguishing examples while building concepts in hippocampal and artificial networks}
\author{}
\date{}

\begin{document}
\maketitle

\begin{abstract}
The hippocampal subfield CA3 is widely regarded as an autoassociative network
that stores experiences as memories, and it receives two anatomically distinct
encodings of every sensory input: a sparse, decorrelated mossy-fiber (MF)
pathway from the dentate gyrus and a dense, correlated perforant-path (PP)
pathway from the entorhinal cortex. Why the circuit would route the same
information along two such different channels has remained unclear, and
previous autoassociative models collapse the two into a single stored pattern
per memory. We propose instead that CA3 stores both MF and PP encodings for
each memory and retrieves either of them depending on the inhibitory tone of
the network, so that sparse MF attractors preserve distinctions between
examples while dense PP attractors merge similar memories into concept-like
representations. We develop a Hopfield-like model with this property, driven
by MF patterns of density 0.02 and PP patterns of density 0.2 and
within-concept correlations of 0.01 and 0.09 respectively. An oscillating
threshold alternates the network between distinct MF examples at high
threshold ($\theta' = 0.5$) and merged PP concepts at low threshold
($\theta' = 0$). We test a direct prediction of this multiplexed code using
public rodent recordings (CRCNS hc-3 and hc-6): across 35 CA3 fields and 99
CA3 W-maze neurons, sparser theta phases convey significantly more
sparsity-corrected information per spike about fine positions and turn
directions, while denser phases convey more information at coarser track
scales, all quantified with two-tailed Wilcoxon signed-rank and Mann-Whitney
U tests. Finally, we import the principle into artificial networks via a
HalfCorr loss that decorrelates only half of the final hidden layer of a
multilayer perceptron; this architecture outperforms both baseline and
fully decorrelated networks on a joint MNIST digit classification and set
identification benchmark, with correlated neurons supporting concept
learning and decorrelated neurons supporting example learning.
\end{abstract}

\section{Introduction}

The capacity of the hippocampus to support vivid recollection of individual
experiences is among the most striking features of mammalian memory
\cite{scoville1957loss}. Within the hippocampal formation, the subfield CA3
has long been regarded as an autoassociative network in which dense recurrent
excitation among pyramidal cells, together with spike-timing-dependent
plasticity at those synapses \cite{bi1998synaptic}, permits the reinstatement
of stored activity patterns from noisy or partial cues
\cite{mcnaughton1987hippocampal,hopfield1982neural}. Sensory information
destined for CA3 enters this circuit along two anatomically segregated routes
from the entorhinal cortex (EC). A direct perforant path (PP) projects from
EC layer II onto the distal dendrites of CA3 pyramidal cells, while an
indirect route passes first through the dentate gyrus (DG) before its mossy
fibers (MF) contact CA3 on more proximal dendrites. The functional
justification for this dual-input architecture, despite decades of study,
has remained unsettled.

The two pathways deliver information with systematically different statistical
properties. DG is equipped with a large population of granule cells and a
high tonic inhibitory tone, which keeps their activity sparse and interference
among overlapping representations low \cite{engin2015tonic,leutgeb2007pattern}.
Sparsification in feedforward projections is in turn associated with
decorrelation, so that MF inputs to CA3 inherit both a sparse and a
decorrelated character. PP connectivity, by contrast, is dense, and natural
correlations between similar sensory inputs should therefore survive the
projection to CA3. Previous theories of the circuit have proposed that the
MF pathway teaches CA3 to store well-separated memory traces during
encoding, while the PP pathway mediates retrieval from a single hybrid
attractor \cite{treves1992computational,mcclelland1995why}. Under these
accounts, each stored experience is represented by one pattern at
equilibrium, irrespective of which route was used to deliver it.

Here we consider an alternative possibility. We ask whether CA3 can store
and retrieve two distinct encodings of the same memory, one MF-like and one
PP-like, and whether the physiological theta oscillation, which rhythmically
modulates inhibitory tone across phases \cite{tsodyks1996population}, serves
as the switch that selects between them. If this were true, sparse,
decorrelated MF encodings and dense, correlated PP encodings would not
collapse into a single stored pattern; instead, the network would multiplex
the two, with different inhibitory thresholds favoring different retrieval
modes. Conceptually, MF attractor basins would remain separate and support
pattern separation suitable for distinguishing examples, while PP basins
would merge along shared features and support concept-like generalization
across similar experiences
\cite{quiroga2005invariant,knowlton1993learning,zeithamova2019brain}. The
theta oscillation, on this view, does not simply pace activity in the
hippocampus; it sweeps the network through a continuum of representational
scales that are all stored at the same synapses.

The paper develops the proposal in three stages. We first construct a
biologically motivated model in which EC, DG, MF and PP encodings are
produced by feedforward networks with realistic statistics, and in which a
Hopfield-like CA3 network \cite{hopfield1982neural,tsodyks1988enhanced}
stores linear combinations of MF and PP patterns. Asynchronous retrieval
dynamics controlled by a threshold \cite{amit1985spin} alternate the network
between MF-like example representations at high threshold and PP-like
concept representations at low threshold. We then test the core prediction
of this account, that single CA3 place cells should convey more information
per spike about fine, example-like positions during sparse theta phases,
and more information about coarse, concept-like positions during dense
phases, using publicly available rodent electrophysiological recordings
\cite{mizuseki2013hc3,karlsson2015hc6}. Finally, we abstract the principle
to artificial neural networks: a DeCorr loss term that decorrelates batch
representations, and a HalfCorr variant that applies the decorrelation to
only half of the final hidden layer, govern a tradeoff between concept-like
and example-like learning in a multilayer perceptron trained on MNIST
\cite{lecun1998gradient}. The contributions of this work are therefore
threefold: an autoassociative model in which a single CA3 network stores
two encodings per memory and alternates between them under threshold
control; empirical confirmation of the model's signature predictions in CA3
place fields across two datasets; and a plug-and-play loss function that
imports the same principle into artificial networks and improves multitask
performance.

\section{Related Work}

The idea that CA3 performs pattern completion through recurrent excitation
traces back to foundational work in attractor networks
\cite{hopfield1982neural,tsodyks1988enhanced,amit1985spin} and to the
complementary learning systems framework
\cite{mcnaughton1987hippocampal,mcclelland1995why}. Several computational
accounts have specifically addressed the role of the MF pathway, typically
concluding that MF inputs support the storage of well-separated patterns that
are later retrieved through PP cues \cite{treves1992computational}. Our
model differs in asserting that two distinct encodings per memory coexist in
CA3 and are retrieved in alternation, rather than collapsing into a single
attractor. Experimental work has in parallel confirmed that DG decorrelates
inputs to CA3 more aggressively than CA3 does on its own, by recruiting
largely non-overlapping granule cell populations when environments differ
\cite{leutgeb2007pattern} and by maintaining the high tonic inhibition that
keeps activity sparse \cite{engin2015tonic}. These findings directly motivate
our choices of pattern statistics for MF and PP inputs in the CA3 model.

Phase precession, in which spikes advance to earlier theta phases as an
animal traverses a place field, is a robust property of hippocampal place
cells \cite{okeefe1993phase,tsodyks1996population} and is accompanied by
phase-dependent changes in tuning sharpness \cite{skaggs1996theta}. Our
analysis extends this literature by partitioning theta into sparse and
dense halves and asking, separately for fine and coarse position binning,
which half of the cycle conveys more sparsity-corrected information. The
prediction tested is the unique signature of a threshold-switched
multiplexed code rather than a generic inhibitory effect. Complementary
evidence from the human medial temporal lobe shows that single neurons can
respond invariantly to different depictions of a particular individual or
concept \cite{quiroga2005invariant}, and behavioral and imaging studies
establish hippocampal involvement in categorization and concept learning
\cite{knowlton1993learning,zeithamova2019brain}. The prototype- and
clustering-based accounts of category learning
\cite{anderson1991adaptive,love2004sustain} provide a natural target
representation for our PP attractors, which merge similar examples along
shared features.

Regularizers that shape hidden representations have been used primarily to
reduce overfitting in deep networks. The DeCov loss decorrelates pairs of
neurons across all inputs \cite{cogswell2016reducing}, whereas the DeCorr
loss we introduce decorrelates pairs of inputs across all neurons in a
layer. The two are complementary rather than equivalent: DeCov improves
generalization on the concept task, while DeCorr induces the opposite
effect by encouraging representations that discriminate examples. Our
HalfCorr variant combines correlated and decorrelated components within a
single layer to resolve the resulting tradeoff, bridging a biological
motivation drawn from CA3 \cite{hopfield1982neural,schapiro2017complementary}
with a concrete machine-learning construction applicable to networks
trained with standard optimizers \cite{kingma2015adam,lecun2015deep}. The
specific interpretation of the theta oscillation we adopt, namely that it
acts as a subtractive threshold on the network, is aligned with the broader
view that inhibitory interneurons shape cortical activity through both
subtractive and divisive effects \cite{isaacson2011inhibition}.

\section{Methods}

\subsection{Transformation of memories along hippocampal pathways}

The memories stored in the model are images from the sneaker, trouser and
coat classes of FashionMNIST \cite{xiao2017fashion}, with 256 images drawn
from each class. These images are first encoded into entorhinal cortex (EC)
patterns by a fully connected linear autoencoder with three hidden layers of
sizes 128, 1024 and 128, batch normalization on each hidden layer, ReLU
nonlinearities on the first and third hidden layers, a sigmoid nonlinearity
on the output layer, and a Heaviside step binarization of the middle hidden
layer with gradients backpropagated through a straight-through estimator
\cite{bengio2013estimating}. The EC layer has $N_{\mathrm{EC}} = 1024$ units
with desired density $a_{\mathrm{EC}} = 0.1$, enforced through a
Kullback-Leibler divergence sparsification term with strength $\lambda = 10$
between the realized and target activation rates. The autoencoder is
optimized for 150 epochs with batch size 64 using Adam
\cite{kingma2015adam} at learning rate $10^{-3}$ and weight decay $10^{-5}$.

The DG, MF and PP encodings are produced from the EC patterns by random
binary connectivity matrices with a winners-take-all activation rule. Each
postsynaptic neuron receives $l$ excitatory synapses of equal strength from
randomly chosen presynaptic neurons; a threshold is then implicitly set so
that the postsynaptic pattern attains the desired density $a_{\mathrm{post}}$.
Taking estimates of rodent hippocampal cell counts as a guide, the DG layer
has $N_{\mathrm{DG}} = 8192$ and the CA3 layer has $N_{\mathrm{CA3}} = 2048$.
We choose densities $a_{\mathrm{DG}} = 0.005$, $a_{\mathrm{MF}} = 0.02$ and
$a_{\mathrm{PP}} = 0.2$, and synapse counts $l_{\mathrm{DG}} = 205$,
$l_{\mathrm{MF}} = 8$ and $l_{\mathrm{PP}} = 205$. We do not directly impose
within-concept correlations; these emerge from propagation and take values
$\rho_{\mathrm{EC}} = 0.15$, $\rho_{\mathrm{DG}} = 0.02$,
$\rho_{\mathrm{MF}} = 0.01$ and $\rho_{\mathrm{PP}} = 0.09$. The relationship
between presynaptic density $a_{\mathrm{pre}}$ and correlation
$\rho_{\mathrm{pre}}$ on one hand and postsynaptic correlation
$\rho_{\mathrm{post}}$ on the other is captured by a closed-form expression
involving the inverse complementary error function (Eq.~1 in the extracted
methods); remarkably, the postsynaptic correlation does not depend on the
size of the network, on the synapse count $l$, or on the absolute threshold
value. Figure~\ref{fig:density_correlation} summarizes these statistics
across regions, and Table~\ref{tab:params} lists their numerical values.

To visualize the activity of the CA3 network back in image space, we train a
separate fully connected feedforward network with one hidden layer of 4096
units to map MF inputs onto their corresponding EC patterns and PP inputs
also onto the same EC patterns, with batch normalization and a ReLU on the
hidden layer and a sigmoid on the output. This visualization network is
optimized for 100 epochs with batch size 128 using Adam at learning rate
$10^{-4}$ and weight decay $10^{-5}$. It is used only to render network
activity as images and is not part of the autoassociative retrieval.

\subsection{Hopfield-like CA3 model}

The CA3 network stores linear combinations
$q^{\mu\nu} = \zeta\, \xi^{\mathrm{MF},\mu\nu} + (1-\zeta)\, \xi^{\mathrm{PP},\mu\nu}$
of MF and PP encodings, with relative strength $\zeta = 0.1$ reflecting the
stronger MF detonator synapses \cite{hopfield1982neural,tsodyks1988enhanced}.
Here $\mu$ indexes concepts and $\nu$ indexes examples. The synaptic matrix
follows the standard Hopfield outer-product form with mean-subtracted
patterns. Cues are formed by randomly flipping the activity of a fraction
$0.01$ of all neurons, and retrieval proceeds by asynchronous Glauber
dynamics \cite{amit1985spin} in which one neuron is chosen per simulation
step and its activation is drawn from a sigmoid whose argument is the
difference between its total recurrent input and a threshold $\theta$,
scaled by an inverse temperature $\beta$. We define rescaled variables
$\theta' = \theta / \left[(1-\zeta)^2 a_{\mathrm{MF}}\right]$ and
$\beta' = \beta \cdot (1-\zeta)^2 a_{\mathrm{MF}}$; simulations use $\beta = 100$
and 10 update cycles, with $\theta' = 0.5$ to retrieve MF patterns and
$\theta' = 0$ to retrieve PP patterns. In the cue-throughout condition,
the cue is applied continuously as external input $h(t) = \sigma x$ with
strength $\sigma = 0.2$ rather than only at initialization.

The overlap between the network activity $S$ and a target pattern $x$ is
$m = (S \cdot x - a\, \sum_i S_i)/(a\, N\, (1-a))$, where $a$ is the density
of the target pattern. To evaluate concept-level retrieval, we define a
target PP concept pattern $\xi^{\mathrm{PP},\mu}$ for concept $\mu$ by
selecting, across all PP examples within $\mu$, the neurons that are most
active until the density $a_{\mathrm{PP}}$ is reached. The theoretical
maximum overlap of the network with $\xi^{\mathrm{PP},\mu}$ is the square
root of the within-concept correlation $\rho_{\mathrm{PP}}$; because this
estimate is asymptotic in the large-network limit, simulations can modestly
exceed it.

\subsection{Experimental data analysis}

Public rodent recordings from two datasets are used to test the model's
predictions. For linear-track analyses we use the CRCNS hc-3 dataset
\cite{mizuseki2013hc3}, and for the W-maze alternation task we use the CRCNS
hc-6 dataset \cite{karlsson2015hc6}. Place cells are selected by requiring
their phase-independent position information per spike, computed using 1-cm
bins across all theta phases, to exceed 0.5. For each selected cell we then
bin spikes into positions (or turn directions) and into theta phases of
width $30^\circ$ for the linear track and $45^\circ$ for the W-maze. To
ensure robust phase-resolved information estimates, we require at least 8
spikes per phase value for the linear-track analyses, at least 5 for the
phase-extracted place fields of Figure~\ref{fig:valid_cells}, and at least 5
per phase value for the W-maze. Theta modulation is required by demanding
that the most active phase contains at least twice as many spikes as the
least active phase.

Information per spike is computed from spike counts $c(r,\phi)$ over spatial
bins $r$ and phase bins $\phi$, together with trajectory occupancies that are
distributed uniformly over $\phi$ \cite{skaggs1996theta}. Because
information estimates are biased by sparsity, we correct them by generating
100 null-matched surrogate neurons in which the spatial bin of each spike is
replaced by a uniformly distributed random value while its phase is
preserved. We subtract the mean null-matched information from the raw
information to obtain a sparsity-corrected estimate. Theta phases are then
partitioned into sparse and dense halves according to the total activity,
the sparsity-corrected information is averaged within each half, a ReLU is
applied to prevent negative values, and the difference between halves is
reported. Shuffled controls permute spike positions and phases to disrupt
any relationship between them, and they confirm that the observed effects
are not an artifact of the binning or of the sparsity correction.

Phase precession is detected by circular-linear regression between spike
progresses and phases, with precession score defined as the mean resultant
length and precession slope defined as the regression slope; cells with
score greater than $0.3$ and slope steeper than $-72^\circ$ per field are
classified as precessing. Statistical comparisons between sparse and dense
halves are carried out with two-tailed Wilcoxon signed-rank tests; the
comparison between precessing and non-precessing fields uses the two-tailed
Mann-Whitney U test, and trends across track scales use Spearman's $\rho$.
All analyses are run separately for the two traveling directions along the
linear track and for inbound and outbound runs along the W-maze center arm.

Bayesian population decoding of W-maze turn direction uses sliding windows
of width $\Delta t = 0.5$ s and stride $0.25$ s. Under the assumption of
independent Poisson-distributed spike counts across neurons and time
\cite{zeithamova2019brain}, the likelihood of each turn direction is
proportional to the standard Poisson population likelihood. We require at
least 2 spikes per window to decode, and we include sessions with at least
5 simultaneously recorded neurons. Before decoding, activity values below
$0.02$ times the maximum for each neuron are rectified to avoid numerical
instabilities. With these criteria, 8 CA3 sessions and 25 CA1 sessions
qualify. Confidence is reported as $|2p(R) - 1|$ and accuracy as 1 when the
maximum-likelihood turn direction matches the ground truth.

\subsection{Artificial neural networks}

The digit classification task uses MNIST \cite{lecun1998gradient} normalized
to zero mean and unit standard deviation across all images and pixels; each
image is additionally assigned a random set number. Set identification, our
example-discrimination task, is evaluated on the training images after
corrupting 20\% of the normalized pixel values to 0; digit classification,
our concept-learning task, is evaluated on the held-out test images.

In the single-task setting, a two-hidden-layer multilayer perceptron with
tanh activations, 50 units per hidden layer and a softmax output is trained
with stochastic gradient descent at batch size 50, learning rate $10^{-4}$
and momentum $0.9$. Training continues until the network reaches greater
than $99.9\%$ accuracy on the training set; as a guide we use 40 epochs for
digit classification, 100 epochs for set identification with 10 sets and
200 epochs for set identification with 50 sets. The loss is the sum of a
cross-entropy term between predicted and true labels and a DeCorr term,
$\mathcal{L} = \mathcal{L}_{\mathrm{CE}} + \lambda\, \mathcal{L}_{\mathrm{DeCorr}}$,
where
\begin{equation}
\mathcal{L}_{\mathrm{DeCorr}} = \frac{1}{B(B-1)} \sum_{\alpha \neq \beta}
\frac{\sum_i (s_{i\alpha} - \bar s_i)(s_{i\beta} - \bar s_j)}
{\sqrt{\left(\sum_i (s_{i\alpha}-\bar s_i)^2 + \epsilon N\right)
\left(\sum_i (s_{i\beta}-\bar s_i)^2 + \epsilon N\right)}},
\end{equation}
the $s_{i\alpha}$ are final hidden-layer activations for image $\alpha$, the
mean $\bar s_i$ is taken across the batch, $\epsilon = 0.001$ aids numerical
convergence and $B$ is the batch size. We use $\lambda = 0$ for the baseline
and $\lambda = 1$ for the DeCorr condition. The DeCov loss of
\cite{cogswell2016reducing} is included for comparison; it differs from
DeCorr by decorrelating pairs of neurons across batch items rather than pairs
of batch items across neurons.

In the multitask setting, the architecture has two fully connected hidden
layers of 100 tanh units and two softmax output heads, one for digit
classification and one for set identification
(Figure~\ref{fig:ml_schematic}). The training dataset contains 1000 images
and 10 sets, and SGD is used at batch size 50 and learning rate $10^{-4}$
for 100 epochs. The loss adds the two cross-entropy heads and either the
DeCorr or the HalfCorr penalty. HalfCorr applies the same DeCorr formula as
above, but restricted to the second half of the final hidden layer, so that
the first 50 units remain correlated and the second 50 units are actively
decorrelated. We use $\lambda = 1$ with tanh activations, $\lambda = 0.04$
with ReLU, $\lambda = 2$ with no nonlinearity and $\lambda = 0$ for the
baseline. Neuron influence on each task is computed as the average decrease
in task accuracy when the activation of that neuron is clamped to zero.

\begin{table}[htbp]
\centering
\caption{Key hippocampus model parameters and their values unless otherwise
noted. Values are those used throughout the model; small deviations in
specific figures are indicated in the relevant results sections.}
\label{tab:params}
\begin{tabular}{lcccc}
\toprule
Region & Size $N$ & Density $a$ & Within-concept corr. $\rho$ & Input synapses $l$ \\
\midrule
EC & 1024 & 0.1 & 0.15 & --- \\
DG & 8192 & 0.005 & 0.02 & 205 \\
CA3 (MF) & 2048 & 0.02 & 0.01 & 8 \\
CA3 (PP) & 2048 & 0.2 & 0.09 & 205 \\
\bottomrule
\end{tabular}
\end{table}

\begin{table}[htbp]
\centering
\caption{Valid cells, fields and sessions passing the inclusion criteria
used in each analysis. All counts are reported as extracted from the
experimental log; CA1 counts serve as comparative controls.}
\label{tab:counts}
\begin{tabular}{lcc}
\toprule
Analysis & CA3 & CA1 \\
\midrule
Phase-extracted place fields (Fig.~\ref{fig:valid_cells}, right) & 35 fields & 47 fields \\
Coarse track scale $1/16$ & 47 cells & 122 cells \\
Coarse track scale $1/8$ & 49 cells & 137 cells \\
Coarse track scale $1/4$ & 56 cells & 144 cells \\
W-maze single-neuron analysis & 99 cells & 187 cells \\
W-maze population decoding (sessions) & 8 & 25 \\
\bottomrule
\end{tabular}
\end{table}

\begin{figure}[htbp]
\centering
\includegraphics[width=0.88\textwidth]{figures/fig_pathways.pdf}
\caption{Schematic of hippocampal pathways to CA3. Entorhinal cortex (EC)
projects directly to CA3 through the perforant path (PP, dense, approximately
4000 synapses per CA3 pyramidal cell), and indirectly via the dentate gyrus
(DG) and its mossy fibers (MF, sparse, approximately 50 granule cell inputs
per CA3 pyramidal cell). Anatomical values are taken from the extraction.}
\label{fig:pathways}
\end{figure}

\begin{figure}[htbp]
\centering
\includegraphics[width=0.78\textwidth]{figures/fig_density_correlation.pdf}
\caption{Encoding statistics across regions used in the model. Densities $a$
and within-concept correlations $\rho$ are those imposed by the
winners-take-all mechanism and measured from the FashionMNIST encodings:
$a_{\mathrm{EC}} = 0.1$, $a_{\mathrm{DG}} = 0.005$, $a_{\mathrm{MF}} = 0.02$,
$a_{\mathrm{PP}} = 0.2$ and $\rho_{\mathrm{EC}} = 0.15$,
$\rho_{\mathrm{DG}} = 0.02$, $\rho_{\mathrm{MF}} = 0.01$,
$\rho_{\mathrm{PP}} = 0.09$. MF encodings are the sparsest and least
correlated; PP encodings are the densest and most correlated.}
\label{fig:density_correlation}
\end{figure}

\section{Results}

The results below follow the three-part program outlined in the
introduction. We first examine the encoding statistics and retrieval
behavior of the CA3 model, then test its predictions in CA3 place-cell
recordings across two datasets, and finally transfer the computational
principle to artificial neural networks. All numerical parameters are
those listed in Table~\ref{tab:params} and the Methods unless explicitly
stated otherwise, and all inclusion counts for the experimental analyses
are listed in Table~\ref{tab:counts}.

\subsection{MF encodings remain distinct while PP encodings build concepts}

We first asked whether the transformations from EC to CA3 assumed in the
model produce MF and PP encodings with the statistics expected from
anatomy. Concretely, we propagated 256 FashionMNIST images from each of
the sneaker, trouser and coat classes through the binary autoencoder to
yield EC patterns, and then through random binary feedforward networks with
winners-take-all activation to yield DG, MF and PP patterns. The resulting
encodings separate cleanly on the two axes of density and within-concept
correlation (Figure~\ref{fig:density_correlation}): MF patterns at CA3 have
density $0.02$ and correlation $0.01$, making them both the sparsest and
the least correlated; PP patterns at CA3 have density $0.2$ and correlation
$0.09$, making them both the densest and the most correlated. A closed-form
relationship derived in the supplementary methods connects the density and
correlation of the presynaptic patterns and the desired density of the
postsynaptic patterns to the correlation of the postsynaptic patterns
through the inverse complementary error function; notably, this
relationship does not depend on network sizes, synaptic density or absolute
threshold values. The empirical statistics obtained from propagating the
images through the network track the theoretical curves across the range
of densities and correlations that we explored.

The key observation is that sparsification in feedforward projections
promotes decorrelation. When we repeat the analysis with randomly generated
presynaptic patterns at various densities $a_{\mathrm{pre}}$ and
correlations $\rho_{\mathrm{pre}}$, we find the same trend as for the
image-derived encodings: decreasing the postsynaptic density monotonically
decreases the postsynaptic correlation. This is consistent with extensive
prior work on biological and artificial feedforward networks
\cite{olshausen1996emergence,leutgeb2007pattern} and motivates the choice
to build the CA3 model around two complementary encodings with the
statistics of MF and PP patterns rather than a single hybrid.

Having obtained MF and PP encodings for each memory, we next stored linear
combinations of the two in a Hopfield-like CA3 network with relative PP
strength $\zeta = 0.1$, and asked whether the network could retrieve either
of them depending on the inhibitory threshold. With MF cues corrupted by
flipping 1\% of the neurons, the network reliably retrieves distinct MF
representations of examples at high threshold ($\theta' = 0.5$) for all
stored loads that we tested, up to 30 cues per concept. At low threshold
($\theta' = 0$), by contrast, the network transitions from retrieving PP
examples at very low loads to retrieving merged PP concepts at moderate
loads: the recovered image increasingly resembles the average image across
all sneakers (or trousers, or coats) rather than any single stored example,
and the overlap with a target PP concept pattern saturates near the
theoretical bound $\sqrt{\rho_{\mathrm{PP}}}$. Using PP cues instead of MF
cues, and using combined cues, yields qualitatively the same pattern of
results, which implies that retrieval in our model depends primarily on the
inhibitory threshold rather than the cue type. This in turn implies that
the network performs heteroassociation between the two encodings of each
memory (Figure~\ref{fig:retrieval_schematic}).

Finally, we replaced the image-derived MF and PP patterns with randomly
generated binary patterns spanning a broader range of statistics, to ensure
that the retrieval behavior is not idiosyncratic to the generation pipeline.
Random MF patterns had density $0.01$ and correlation $0$, while PP
patterns had density $0.5$. In these networks, MF and PP examples both
exhibit a characteristic capacity above which retrieval fails, while PP
concepts exhibit a characteristic minimum number of stored examples below
which the concept cannot be built. At intermediate loads, there exists a
multiscale retrieval regime in which both MF examples and PP concepts are
recoverable, and this regime widens as the PP correlation increases.

\begin{figure}[htbp]
\centering
\includegraphics[width=0.95\textwidth]{figures/fig_retrieval_schematic.pdf}
\caption{Threshold control of retrieval in the Hopfield-like CA3 model.
Left: a theta-like oscillation in the rescaled activity threshold alternates
between a high value ($\theta' = 0.5$) that favors MF example retrieval and
a low value ($\theta' = 0$) that favors PP concept retrieval. Right: at
high threshold the network retrieves the sparse, decorrelated MF
representation of the cued memory; at low threshold, a dense, correlated
PP representation emerges that merges examples along shared features.}
\label{fig:retrieval_schematic}
\end{figure}

\subsection{An oscillating threshold alternates between examples and concepts}

The previous section established that the same CA3 network can be driven
into either MF-like or PP-like retrieval depending on the threshold. We
next asked whether a physiologically motivated, time-varying threshold can
make the network alternate between the two retrieval modes within a single
simulation, which would make the account testable against phase-dependent
physiology. To this end we ran the retrieval dynamics while varying the
threshold through four scenarios: a baseline condition in which the
threshold switched abruptly between $0.5$ (sparse/MF) and $0$ (dense/PP),
a sinusoidal variant with the same extrema, a variant with narrower
extrema of $0.55$ and $0.25$, and a cue-throughout variant in which a
weakened MF cue continued to drive the network throughout the simulation
rather than only at initialization. In each scenario, 20 cues were tested
in each of 20 networks with 50 randomly chosen examples stored from each
of 3 concepts, for a total of 150 examples per network.

Under the baseline and sinusoidal conditions, the network alternates
between MF example and PP concept representations on the same schedule as
the threshold. During the low-threshold phase, the overlap with a target
PP concept approaches the theoretical bound, and visualization of the
retrieved pattern reveals a concept-like prototype. During the high-threshold
phase, the overlap instead concentrates on individual MF examples, and from
one oscillation cycle to the next the retrieved MF pattern can hop among
different examples of the same concept: the low-threshold phase preserves
the concept identity of the cue but loses its example identity, so
different examples are equally likely at the subsequent high-threshold
phase. When the MF cue is applied throughout the simulation, this
multistability disappears and the network settles on the target example;
when the oscillation amplitude is reduced (the $0.55$/$0.25$ scenario),
alternation is disrupted and the network instead stalls on one encoding
type. These behaviors were summarized over update cycles 60 to 120, well
after initialization transients had decayed.

These simulations yield three experimentally testable predictions for CA3
physiology. First, sparse theta phases should preferentially encode fine,
example-like features of an experience, while dense phases should
preferentially encode coarse, concept-like features. Second, the transition
between scales should be driven by the inhibitory threshold in a
subtractive rather than a divisive manner, because divisive modulation
would scale tuning curves without altering their sharpness. Third, lower
oscillation amplitudes should reduce access to the range of representational
scales, consistent with reports that stronger theta power is associated
with better memory retrieval \cite{okeefe1993phase,tsodyks1996population}.

\subsection{CA3 place cells encode fine positions during sparse theta phases}

To test the first prediction, we turned to publicly available rodent CA3
place-cell recordings from linear-track navigation
\cite{mizuseki2013hc3,skaggs1996theta}. We construe each CA3 neuron as
storing spatial memories that, like FashionMNIST images, can be organized
at different scales: fine positions within a place field play the role of
examples, while coarse regions of the track play the role of concepts. The
model therefore predicts that, within a place field, sparse theta phases
convey more sparsity-corrected position information per spike than dense
phases. A priori this is a nontrivial prediction: if theta acted divisively,
sharper tuning might be expected during denser phases instead
\cite{isaacson2011inhibition}, as seen for visual tuning curves under
attention.

We selected place cells with phase-independent position information above
$0.5$ bits per spike using 1-cm bins, and extracted 35 valid CA3 fields and
47 valid CA1 fields that passed both this criterion and the theta-modulation
requirement (Figure~\ref{fig:valid_cells}, right). For each field, we
partitioned theta into sparse and dense halves, computed sparsity-corrected
position information per spike within each half, and tested the difference
with a two-tailed Wilcoxon signed-rank test. CA3 place fields convey
significantly more position information during sparse theta phases than
during dense phases, and this effect is present for both phase-precessing
fields and non-precessing fields, although the effect is slightly
non-significant in the latter. Precession was detected by circular-linear
regression with thresholds of score greater than $0.3$ and slope steeper
than $-72^\circ$ per field. The difference between precessing and
non-precessing fields was assessed with the two-tailed Mann-Whitney U test.
Critically, the same effect is absent in data in which spike phases and
positions are shuffled, which rules out a sparsity-correction artifact,
and it is also absent in CA1 place fields, which shows that the signature
is specific to CA3 and is not simply inherited through projection.

When we next asked about coarse positions, partitioning the whole track
into 4 contiguous bins at track scales $1/16$, $1/8$ and $1/4$, we recovered
47, 49 and 56 valid CA3 place cells and 122, 137 and 144 valid CA1 place
cells respectively at those scales (Figure~\ref{fig:valid_cells}, left). At
the coarsest scales, dense theta phases convey more sparsity-corrected
position information per spike than sparse phases, exactly reversing the
trend seen at fine scales. The trend across scales is monotonic, with a
significant Spearman's $\rho$, and is absent from shuffled controls. CA1
place cells show the same dense-phase preference at coarse scales. The
emergent picture is that the same CA3 neurons represent fine,
example-like features during sparse theta phases and coarse, concept-like
features during dense phases, consistent with the model's multiplexed
code.

\begin{figure}[htbp]
\centering
\includegraphics[width=0.95\textwidth]{figures/fig_valid_cells.pdf}
\caption{Numbers of valid cells and fields passing the inclusion criteria
described in Methods. Left: coarse-position analysis at track scales
$1/16$, $1/8$, $1/4$ yields 47, 49 and 56 valid CA3 place cells and 122,
137 and 144 valid CA1 place cells, respectively. Right: the fine-scale
place field analysis yields 35 valid CA3 fields and 47 valid CA1 fields,
while the W-maze task yields 99 valid CA3 neurons and 187 valid CA1
neurons.}
\label{fig:valid_cells}
\end{figure}

\subsection{W-maze turn direction is preferentially encoded in sparse phases}

The linear-track results support our model in the single modality of
position. To test whether the prediction generalizes to a non-spatial
feature, we analyzed CA3 place cells recorded during a W-maze alternation
task in which mice must alternately visit the left and right arms between
runs along the center arm \cite{karlsson2015hc6}. In this task,
center-arm place cells are known to encode the turn direction in addition
to position, so turn direction plays the role of a fine example feature
that distinguishes otherwise similar runs, and position along the center
arm plays the role of a coarse concept feature shared across turn
directions. The model therefore predicts that sparse theta phases will
preferentially encode turn direction, while dense phases will represent
position without regard to the upcoming (or past) turn.

We partitioned spikes into 2 turn directions and theta-phase bins of
width $45^\circ$, requiring at least 5 spikes per phase value and at least
twice as many spikes at the most active phase as at the least active phase.
These criteria yield 99 valid CA3 neurons and 187 valid CA1 neurons for
the single-neuron analysis. For each cell, we computed turn information
per spike at each phase, corrected it with 100 null-matched samples that
randomize the turn direction of each spike while preserving its phase, and
averaged the sparsity-corrected turn information across the sparse and the
dense halves of the theta cycle. Across both outbound and inbound runs,
CA3 neurons convey significantly more turn information per spike during
sparse phases than during dense phases, as assessed by two-tailed Wilcoxon
signed-rank tests. The comparison between outbound and inbound runs is
carried out with the Mann-Whitney U test. The same effect is not present
in shuffled controls, and the difference between sparse and dense phases
is significantly larger in CA3 than in CA1, consistent with the idea that
the predicted signature emerges in CA3 and is not trivially propagated
forward.

At the population level, we asked whether this single-neuron effect
translates into a more confident and more accurate Bayesian decoding of
turn direction from the sparse half of the cycle. We included sessions
with at least 5 simultaneously recorded neurons, yielding 8 valid CA3
sessions and 25 valid CA1 sessions, and decoded in sliding windows of
width $0.5$ s and stride $0.25$ s, requiring at least 2 spikes per window.
Under the Poisson population-likelihood model, the CA3 decoder exhibits
greater confidence about turn direction during sparse theta phases than
during dense phases, and its maximum-likelihood estimate is more often
correct during sparse phases. Both effects are specific to CA3: the CA1
population does not show the same pattern, consistent with the subfield
specificity seen at the single-neuron level. Taken together, the linear
track and W-maze analyses provide model-specified, dataset-independent
evidence for the multiplexed code across two tasks, two encoding
modalities, and two levels of spatial resolution.

\subsection{Complementary encodings improve multitask machine learning}

Having found the signatures of the model in CA3, we asked whether the same
computational principle can improve learning in artificial networks, far
removed from hippocampal anatomy. We augmented MNIST with a random set
number for each image and asked a multilayer perceptron to perform either
digit classification, a concept task that requires building generalizations
over many examples, or set identification, an example task that requires
distinguishing images that may be otherwise visually similar
(Figure~\ref{fig:ml_schematic}). Concept performance was measured on
held-out test images, while example performance was measured on training
images with 20\% of their normalized pixel values randomly set to 0, which
forces the network to commit to an example-level representation rather
than to rely on memorization of exact pixel values.

In the single-task setting, with two tanh hidden layers of 50 units each,
the baseline network with $\lambda = 0$ performs better on digit
classification (the concept task), while the DeCorr network with
$\lambda = 1$ performs better on set identification (the example task).
The effect varies consistently with $\lambda$: larger DeCorr strength
further improves example learning and further degrades concept learning.
DeCov, which decorrelates pairs of neurons across inputs rather than pairs
of inputs across neurons, improves digit classification, as expected of a
generalization-promoting regularizer, but does not substantially improve
set identification, confirming that DeCorr and DeCov are distinct
interventions. Each reported condition averages over 32 independent
networks trained with SGD at batch size 50, learning rate $10^{-4}$ and
momentum $0.9$ to greater than $99.9\%$ training accuracy.

In the multitask setting, with two tanh hidden layers of 100 units each
and a joint dataset of 1000 training images labeled with both digit and
one of 10 sets, the tradeoff becomes explicit. A DeCorr network is better
than baseline on the example task but worse on the concept task, as
expected. A HalfCorr network, in which the DeCorr penalty is applied only
to the second 50 units of the final hidden layer, removes this tradeoff:
it outperforms both baseline and DeCorr on the example task and matches or
exceeds baseline on the concept task, over 64 networks per condition. The
advantage is preserved across a range of architectural choices, including
ReLU hidden layers at $\lambda = 0.04$ and purely linear hidden layers at
$\lambda = 2$, which suggests that HalfCorr is not tuned to a particular
activation function but rather captures a more general computational
principle.

The mechanism by which HalfCorr achieves this is visible at the
single-neuron level. When we clamp each neuron's activation to zero in
turn and measure the decrease in accuracy on each task, we find that the
correlated half of the final hidden layer is more influential on digit
classification, while the decorrelated half is more influential on set
identification. In other words, the network spontaneously learns to use
its two subpopulations of neurons for the two tasks for which their
correlation statistics are best suited, in direct parallel with the
differential contributions of MF and PP encodings in CA3. We did not
separately manipulate the sparsities of the two halves; in some
configurations the decorrelated half is also sparser than the correlated
half, but the primary control is correlation rather than sparsity.

\begin{figure}[htbp]
\centering
\includegraphics[width=0.90\textwidth]{figures/fig_ml_schematic.pdf}
\caption{Multitask multilayer perceptron architecture. Two fully connected
hidden layers of 100 tanh units feed two softmax output heads, one for
digit classification and one for set identification. The HalfCorr loss
applies the DeCorr penalty only to the second half of the final hidden
layer, yielding a correlated half (PP-like, orange) and a decorrelated
half (MF-like, purple). The DeCorr baseline applies the penalty to all
neurons, while the baseline condition applies no penalty.}
\label{fig:ml_schematic}
\end{figure}

\section{Discussion}

Our central proposal is that the dual MF and PP inputs to CA3 support a
multiplexed code in which each memory is stored twice, once as a sparse,
decorrelated MF pattern and once as a dense, correlated PP pattern, and
retrieved either as an example or as a concept depending on the inhibitory
tone of the network. At the level of the circuit, this reframes the theta
oscillation not as a clock that gates a single stored pattern per memory
but as a sweep across representational scales, with sparser phases
preserving example identity and denser phases preserving concept identity.
At the level of behavior, it predicts that the hippocampus should be able
to alternate between barcode-like episodic retrieval and prototype-like
categorical retrieval without modifying its synapses, simply by tracking
changes in inhibitory tone. The place-cell analyses reported here are
consistent with this reframing across two datasets, two behavioral tasks
and two levels of spatial resolution, and the HalfCorr results show that
the same principle can be transplanted into artificial networks without
the full mechanistic apparatus of recurrent attractor dynamics.

Our work connects a specific anatomical feature of the hippocampal circuit
to foundational attractor theory
\cite{hopfield1982neural,tsodyks1988enhanced,amit1985spin}. Prior
computational models have either assigned MF and PP inputs to different
operating regimes, typically MF for storage and PP for retrieval from a
single hybrid attractor \cite{treves1992computational}, or have proposed
that both example-like and concept-like representations coexist but
converge at CA1 rather than at CA3 \cite{schapiro2017complementary}. Our
model diverges from both of these by positing that two encodings of each
memory coexist at CA3 itself, are retrieved by the same recurrent network
under different thresholds, and manifest as the sparse- and dense-phase
information signatures that we identify in CA3 place cells. Because the
connection to statistical-mechanics accounts of capacity in sparse networks
is explicit, our model also inherits their predictions for how capacity
scales with network size and pattern density, making it possible in
principle to calibrate it against large-scale recordings of CA3 activity.

The analyses of theta phase coding reported here enrich the existing
literature on phase precession and phase-dependent tuning in several ways
\cite{okeefe1993phase,tsodyks1996population,skaggs1996theta}. Previous
work has largely focused on CA1, and the subfield-specific signatures we
report in CA3, including the dense-phase preference for coarse positions
and the sparse-phase preference for turn direction in the W-maze, are not
obviously inherited from CA1 and are not predicted by accounts of theta
that treat it purely as a divisive modulator \cite{isaacson2011inhibition}.
The observation that some CA3 neurons encode turn direction only at
certain theta phases, rather than across the whole cycle, suggests that
the population of cells traditionally labeled as splitter cells may be
broader than is apparent from phase-averaged analyses, because neurons
that are theta-specific splitters will be discarded by whole-cycle
analyses. This resonates with broader views of the hippocampus as
representing the external world in a richly multiplexed fashion across
oscillation phases.

The extension to artificial neural networks, mediated by the DeCorr and
HalfCorr loss terms, deserves discussion in its own right. DeCorr and the
previously studied DeCov loss \cite{cogswell2016reducing} operate on
different axes of the activation matrix, and the contrast is empirically
important: DeCov reduces overfitting by encouraging different neurons to
respond to different inputs, while DeCorr actively pulls apart the
representations of different inputs within the same layer, which is
useful when the task is to tell those inputs apart. HalfCorr, which
applies DeCorr to only a designated subset of neurons, realizes within a
single layer the kind of functional heterogeneity that is more commonly
achieved by modular architectures \cite{lecun2015deep}. It is natural to
speculate that a spectrum of decorrelation strengths, assigned neuron by
neuron, could produce a continuous family of representational scales
analogous to the continuum swept out by a theta-like inhibitory
oscillation in CA3. Whether this heterogeneity would be useful in
continual learning, reinforcement learning or natural biological learning
is beyond the scope of the current paper, but the results reported here
suggest that correlation-based regularizers are an underexplored design
dimension for multitask learning.

Several important limitations of the present work should be noted. Our
model does not include CA1 in its circuitry, and although we report
comparative analyses in this subfield that serve as controls, a fuller
account of how CA3 outputs are combined with temporoammonic inputs at
CA1 is needed to connect the proposal to behavior and to downstream
neocortical targets. Our model also treats memory storage and retrieval
as separate simulation regimes, which is a strong simplification given
that storage and retrieval likely coexist at different theta phases in
vivo \cite{tsodyks1996population}. We have not attempted to fit the
model's parameters to any single empirical dataset, preferring instead to
choose values that reflect the anatomy in order of magnitude and to
verify that the qualitative phenomena are preserved across a wide range.
The machine-learning experiments likewise use relatively small MNIST
subsets and a simple multilayer perceptron, chosen to expose the
representational principle rather than to compete for state-of-the-art
accuracy. Scaling HalfCorr to larger, more naturalistic tasks, and to
architectures with depth and attention, is a natural next step.

Taken together, the biological and computational results reported here
argue that multiplexing two encodings of each memory at CA3, and using the
theta oscillation to select between them, provides a principled account of
why CA3 receives two systematically different inputs and how it can
support both example- and concept-level retrieval. The success of HalfCorr
as a plug-and-play loss function further suggests that the same principle
can benefit artificial systems that must simultaneously distinguish similar
inputs and generalize over them \cite{quiroga2005invariant,lecun2015deep}.
More broadly, the results support the view that cortical inhibition should
not be treated simply as a damper on activity but as an active, oscillatory
control variable that selects among coexisting representational scales
\cite{isaacson2011inhibition}. Future work combining large-scale recordings
of CA3 during theta, fits of the Hopfield-like model to those recordings,
and scaled-up HalfCorr-like losses in large neural networks should make
it possible to sharpen and test this view.

\bibliography{references}

\end{document}
